% =============================================================================
% METHODOLOGY - CITA Paper in ALKALI Format
% =============================================================================

\vspace{-2mm}
\section{Methodology}
\label{sec:methodology}

We present the CITA training methodology, organized into three phases: (1) setup and data preparation, (2) instruction-tuned formatting, and (3) contrastive training configuration.

\vspace{-2mm}
\subsection{Phase 0: Model and Infrastructure Setup}
\label{sec:phase0}

\begin{table}[H]
\centering
\small
\setlength{\tabcolsep}{6pt}
\renewcommand{\arraystretch}{1.2}
\begin{tabular}{@{}ll@{}}
\toprule
\textbf{Component} & \textbf{Specification} \\
\midrule
Base Model & Llama-3.1-8B~\cite{dubey2024llama} \\ \hline
Prompt Format & Chat-style (system/user/assistant) \\ \hline
Tokenizer & LLaMA-3 with BOS/EOS tokens \\ \hline
Hardware & NVIDIA A100-40GB / A100-80GB \\ \hline
Training Data & PKU-SafeRLHF~\cite{ji2024pku} \\
\bottomrule
\end{tabular}
\caption{Training infrastructure and configuration.}
\label{tab:infrastructure}
\vspace{-4mm}
\end{table}

\textbf{Rationale.} We select Llama-3.1-8B for its strong instruction-following capabilities and accessibility, following the Llama model family tradition~\cite{touvron2023llama,touvron2023llama1,dubey2024llama3}. PKU-SafeRLHF provides high-quality preference pairs with explicit safety annotations~\cite{bai2022training}, enabling evaluation of instruction-conditioned safety behaviors. Alternative datasets like OpenAssistant~\cite{kopf2023openassistant}, UltraFeedback~\cite{cui2023ultrafeedback}, and Tulu~\cite{ivison2023camels} could also be used.

\vspace{-2mm}
\subsection{Phase 1: Dataset Formatting}
\label{sec:phase1}

We define two complementary data formats that enable instruction-conditioned training:

\paragraph{Instruction-Tuned Alignment (ITA) Format.}
Each training example explicitly includes an alignment instruction alongside the user query:

\begin{small}
\begin{verbatim}
<user>
### Alignment Instruction:
${alignment_instruction}
### User Prompt:
${user_prompt}
<assistant>
${expected_response}
\end{verbatim}
\end{small}

\paragraph{Contrastive ITA Format.}
For preference learning, we extend ITA to include rejected responses:

\begin{small}
\begin{verbatim}
### Rejected Response (Bad):
${rejected_response}
### Expected Response (Good):
${accepted_response}
\end{verbatim}
\end{small}

\paragraph{Example-Based Alignment (EBA) Format.}
Standard DPO format without explicit instructions for baseline comparison:
\begin{center}
\texttt{\{prompt, chosen, rejected\}}
\end{center}

\vspace{-2mm}
\subsection{Phase 2: Training Configuration}
\label{sec:phase2}

\begin{table}[H]
\centering
\scriptsize
\setlength{\tabcolsep}{4pt}
\renewcommand{\arraystretch}{1.2}
\begin{tabularx}{\columnwidth}{@{}l*{5}{>{\centering\arraybackslash}X}@{}}
\toprule
\textbf{Parameter} & \textbf{SFT} & \textbf{PPO} & \textbf{GRPO} & \textbf{DPO} & \textbf{CITA} \\
\midrule
Learning Rate & 2e-4 & 1e-5 & 5e-6 & 1e-5 & 5.4e-6 \\ \hline
Batch $\times$ Accum & 2$\times$4 & 16$\times$4 & 12$\times$2 & 1$\times$8 & 1$\times$8 \\ \hline
Epochs & 1 & 1 & 1 & 1 & 1 \\ \hline
LoRA Rank & 16 & 16 & 16 & 16 & 16 \\ \hline
$\beta$ (temp.) & --- & --- & --- & 0.1 & 0.107 \\ \hline
KL Coef. & --- & 0.1 & --- & --- & 2.3e-4 \\
\bottomrule
\end{tabularx}
\caption{Training hyperparameters for all 5 alignment methods. All use LoRA (r=16, $\alpha$=16) for parameter-efficient fine-tuning.}
\label{tab:training_config}
\vspace{-4mm}
\end{table}

\textbf{Key Design Choice.} The CITA stage uses significantly lower learning rates (3--4$\times$ lower than typical DPO) because instruction-augmented sequences are 30--40\% longer, producing larger gradients. This finding emerged from hyperparameter optimization using Optuna~\cite{akiba2019optuna} (Section~\ref{sec:hyperparams}). We use parameter-efficient fine-tuning via LoRA~\cite{hu2022lora,dettmers2023qlora} rather than full fine-tuning to enable single-GPU training. Alternative PEFT methods include adapter tuning~\cite{houlsby2019parameter}, prompt tuning~\cite{lester2021power}, and prefix tuning~\cite{li2021prefix}.
